\documentclass[11pt]{amsart}

\usepackage[margin=1in]{geometry}
\usepackage{amsmath,amssymb}
\usepackage{booktabs}
\usepackage{array}
\usepackage[colorlinks=true,allcolors=blue]{hyperref}

\emergencystretch=3em
\tolerance=1200
\hbadness=4000

\newcommand{\TT}[1]{\mathrm{TT}_{#1}}
\newcommand{\Z}{\mathbb{Z}}
\DeclareMathOperator{\Cay}{Cay}

\theoremstyle{plain}
\newtheorem{theorem}{Theorem}[section]
\newtheorem{lemma}[theorem]{Lemma}
\newtheorem{proposition}[theorem]{Proposition}
\theoremstyle{definition}
\newtheorem{definition}[theorem]{Definition}
\theoremstyle{remark}
\newtheorem{remark}[theorem]{Remark}
\newtheorem{question}[theorem]{Question}

\begin{document}

\title[An explicit oriented graph gives $k(3,6)\geq 73$]{An explicit oriented graph gives
$k(3,6)\geq 73$ (Erd\H{o}s Problem \#112)}

\author{Daniel Kirtchakov}
\address{Independent researcher (\texttt{05oz}); no institutional affiliation. ORCID:
\href{https://orcid.org/0009-0009-5213-4098}{0009-0009-5213-4098}.}
\email{daniel@halfounce.io}
\urladdr{https://halfounce.io}

\thanks{\emph{Certify Part O, v0.19.0.} Each graph is checked by two independent checkers on one
machine; a replay on a second machine has not been made. The certificates and scripts are in the
folder \texttt{k36/} of \url{https://github.com/05oz/certify}, archived on Zenodo (all versions:
\href{https://doi.org/10.5281/zenodo.21799111}{doi:10.5281/zenodo.21799111}).}

\thanks{\emph{Computation and authorship.} The graphs, the scripts in \texttt{scripts/} and this
note were produced with Claude (Anthropic), directed by the author. How the graphs were found plays
no part in their validity and is not described here: each statement is an explicit graph in a file
that the checkers of Section~\ref{sec:verify} accept.}

\begin{abstract}
Let $k(n,m)$ be the least $k$ such that every oriented graph on $k$ vertices contains $n$ pairwise
non-adjacent vertices or a transitive tournament on $m$ vertices (Erd\H{o}s Problem \#112). We give
a Cayley digraph $G_{72}$ on $\Z_{24}\times\Z_3$ with $25$ generators that has no three pairwise
non-adjacent vertices and no transitive tournament on six vertices, so $k(3,6)\geq 73$; the largest
published lower bound found is $71$. The first property has a short proof by hand. The second is
checked on the whole graph by two independently written programs; by vertex-transitivity it
amounts to a check of a $25$-vertex subdigraph. We also record two graphs that show
$k(7,3)\geq 39$ and $k(8,3)\geq 47$, bounds first published by T.~S.~Orden, with explicit
isomorphisms to his graphs.
\end{abstract}

\maketitle

\noindent\textbf{Question.} Erd\H{o}s Problem \#112 asks for the least number $k(n,m)$ of vertices
that forces every directed graph to contain $n$ independent vertices or a transitive tournament on
$m$ vertices. Exact values are known only for a few small pairs. This note gives an explicit graph
for one lower bound, $k(3,6)\geq 73$. It also contains two graphs for $k(7,3)\geq 39$ and
$k(8,3)\geq 47$. Those two bounds are T.~S.~Orden's, published on 2026-11-20 \cite{Orden2};
Section~\ref{sec:credit} says how his graphs and these are related.

\section{The number \texorpdfstring{$k(n,m)$}{k(n,m)}}\label{sec:def}

The problem page \cite{Blo} (page and forum thread read in full on 2026-11-11) defines
$k=k(n,m)$ as minimal such that any directed graph on $k$ vertices contains an independent set of
size $n$ or a transitive tournament of size $m$, asks for its value, and attributes the problem to
Erd\H{o}s and Rado \cite{ErRa67}. Ihringer, Rajendraprasad and Weinert \cite{IRW21} write the same
number as $r(I_m,L_n)$, with the size of the independent set written before the size of the
tournament: the least order at which every oriented graph contains an independent set of
cardinality $m$ or a transitive tournament on $n$ vertices. Oriented means no loops and at most one
arc between two vertices; a footnote in \cite{IRW21} gives this as the reason for writing
``oriented'' rather than ``directed''. So $k(n,m)=r(I_n,L_m)$, and \cite{Orden}, \cite{Orden2} and
\cite{rsv} (Section~\ref{sec:credit}) use the same definition.

A set of vertices is independent when no arc joins two of its members; a transitive tournament
$\TT{m}$ is a set of $m$ vertices ordered $v_1,\dots,v_m$ with $v_i\to v_j$ whenever $i<j$. An
oriented graph on $N$ vertices with neither an independent $n$-set nor a $\TT{m}$ shows
$k(n,m)\geq N+1$. The graphs below are oriented, so each is a witness under either wording.

\section{Result}\label{sec:result}

\begin{theorem}\label{thm:main}
$k(3,6)\geq 73$.
\end{theorem}

The proof is an explicit oriented graph, $G_{72}$ (Section~\ref{sec:graphs}). Two more graphs,
$G_{38}$ and $G_{46}$, show $k(7,3)\geq 39$ and $k(8,3)\geq 47$. Orden published those two bounds
before this release \cite{Orden2}, and the two graphs are his graphs up to relabelling
(Section~\ref{sec:credit}); they are in this release as a record, not as a claim. Below, $\alpha$
is the independence number and $\tau$ the number of vertices of the largest transitive
subtournament; both are exact values computed by \texttt{verify\_bounds.py}. The last two columns
are the published values of Section~\ref{sec:credit}.

\begin{center}
\footnotesize
\begin{tabular}{@{}lrrrrl>{\raggedright\arraybackslash}p{2.3cm}>{\raggedright\arraybackslash}p{3.6cm}@{}}
\toprule
graph & vertices & arcs & $\alpha$ & $\tau$ & bound & largest published lower bound found & published upper bound\\
\midrule
$G_{72}$ & 72 & 1800 & 2 & 5 & $k(3,6)\geq 73$ & 71 \cite{Orden2} & 161 \cite{IRW21}; 137 with $k(3,4)=21$ \cite{PartG}\\
$G_{38}$ & 38 & 228 & 6 & 2 & $k(7,3)\geq 39$ & 39 \cite{Orden2} & 45 \cite{IRW21}\\
$G_{46}$ & 46 & 322 & 7 & 2 & $k(8,3)\geq 47$ & 47 \cite{Orden2} & 59 \cite{IRW21}\\
\bottomrule
\end{tabular}
\end{center}

\noindent The files are \texttt{certificates/k36\_N72.json}, \texttt{certificates/k73\_N38.json} and
\texttt{certificates/k83\_N46.json}. $\tau=2$ means that $G_{38}$ and $G_{46}$ have no transitive
triangle: every triangle in them is a directed $3$-cycle.

\section{The three graphs}\label{sec:graphs}

\begin{definition}[$G_{38}$, over $\Z_{19}$]
Vertices are the pairs $(i,a)$ with $i\in\Z_{19}$ and $a\in\{0,1\}$, numbered $19a+i$ in the file.
There is an arc $(i,a)\to(i+d \bmod 19,\,b)$ exactly when $d\in S_{ab}$, where
$S_{00}=\{8,12,13,17\}$, $S_{01}=\{2,6\}$, $S_{10}=\{1,5\}$ and $S_{11}=\{8,12,14,18\}$. Every vertex
has out-degree $6$ and in-degree $6$, and $(i,a)\mapsto(i+1,a)$ is an automorphism with two orbits.
\end{definition}

\begin{definition}[$G_{46}$, over $\Z_{23}$]
The same rule with $\Z_{23}$ and numbering $23a+i$, where $S_{00}=\{1,3,11,18\}$,
$S_{01}=S_{10}=\{6,14,16\}$ and $S_{11}=\{1,3,11,16\}$. Every vertex has out-degree $7$ and
in-degree $7$.
\end{definition}

\begin{definition}[$G_{72}$, a Cayley digraph on $G=\Z_{24}\times\Z_3$]
The element $(x,y)$ is vertex $3x+y$ in the file, and there is an arc $g\to g+s$ for every $s$ in
\begin{multline*}
S=\{(0,2),(2,0),(2,1),(2,2),(3,0),(5,1),(6,1),(6,2),(8,0),(8,1),(9,2),(10,0),(10,1),\\
(11,0),(11,1),(11,2),(12,2),(14,1),(15,2),(16,1),(18,0),(19,0),(19,1),(21,1),(21,2)\}.
\end{multline*}
So $|S|=25$, and every vertex has out-degree $25$ and in-degree $25$.
\end{definition}

\subsection*{Checks by hand}

\begin{lemma}[Orientation]\label{lem:orient}
The three digraphs have no loops and no $2$-cycles.
\end{lemma}

\begin{proof}
$0$ is in neither $S_{00}$ nor $S_{11}$, and $S_{00}\cap-S_{00}$, $S_{11}\cap-S_{11}$ and
$S_{01}\cap-S_{10}$ are empty; $(0,0)\notin S$ and $S\cap-S$ is empty. So there are no loops and no
$2$-cycles.
\end{proof}

\begin{lemma}[No transitive triangle in $G_{38}$ and $G_{46}$]\label{lem:tt3}
A transitive triangle $u\to v\to w$, $u\to w$ with $u,v,w$ in orbits $a,b,c$ translates to
$u=(0,a)$, $v=(d,b)$, $w=(d+e,c)$ with $d\in S_{ab}$, $e\in S_{bc}$ and $d+e\in S_{ac}$, and any such
$d$ and $e$ give one. So there is none exactly when the eight sets $(S_{ab}+S_{bc})\cap S_{ac}$,
taken mod $q$, are empty: a short hand computation.
\end{lemma}

\begin{lemma}[No three independent vertices in $G_{72}$]\label{lem:alpha}
$\alpha(G_{72})=2$.
\end{lemma}

\begin{proof}
Vertices $g\neq h$ are non-adjacent exactly when $h-g$ lies in
$T=G\setminus(S\cup-S\cup\{(0,0)\})$, so three independent vertices exist exactly when $T$ contains
$t_1$, $t_2$ and $t_1+t_2$. Here
$T=(\{1,4,7,17,20,23\}\times\Z_3)\cup\{(9,0),(12,0),(15,0)\}$, with $21$ elements. The $\Z_{24}$
coordinates of $T$ lie in $F=\{1,4,7,9,12,15,17,20,23\}$, and
$F+F=\{0,2,3,5,6,8,10,11,13,14,16,18,19,21,22\}$ does not meet $F$. So $\alpha(G_{72})\leq 2$; and
$T$ is not empty, so $\alpha(G_{72})=2$.
\end{proof}

\begin{lemma}[Reduction for transitive tournaments]\label{lem:reduce}
Let $D[S]$ be the subdigraph of $G_{72}$ induced on $S$, the out-neighbourhood of $(0,0)$. Then
$\tau(G_{72})=\tau(D[S])+1$. In particular $G_{72}$ contains a $\TT{6}$ exactly when $D[S]$, which
has $25$ vertices, contains a $\TT{5}$.
\end{lemma}

\begin{proof}
$G_{72}$ is a Cayley digraph, so each translation $g\mapsto g+h$ is an automorphism. A transitive
tournament has exactly one vertex with no in-neighbour in it, its source. Translating the source
to $(0,0)$ puts the other vertices in $S$, where they form a transitive tournament with one vertex
fewer. Conversely, $(0,0)$ has an arc to every vertex of $S$, so a transitive tournament in $D[S]$
together with $(0,0)$ is a transitive tournament with one vertex more.
\end{proof}

\noindent\emph{The computer part} is what remains: $\alpha(G_{38})\leq 6$, $\alpha(G_{46})\leq 7$,
and no $\TT{5}$ in $D[S]$ (equivalently, no $\TT{6}$ in $G_{72}$). Both checkers of
Section~\ref{sec:verify} test the whole graph $G_{72}$; \texttt{verify\_bounds.py} finds
$\tau(G_{72})=5$, so $\tau(D[S])=4$ by Lemma~\ref{lem:reduce}.

\begin{proof}[Proof of Theorem~\ref{thm:main}]
By Lemmas~\ref{lem:orient} and~\ref{lem:alpha}, $G_{72}$ is an oriented graph on $72$ vertices
with no three pairwise non-adjacent vertices, and by the computation of Section~\ref{sec:verify}
it has no transitive tournament on six vertices. So $k(3,6)\geq 72+1=73$.
\end{proof}

\noindent\texttt{scripts/check\_constructions.py} rebuilds the three arc sets from the data above
(written into the script, not read from the files), checks that they equal the arc lists in the
files, and checks Lemmas~\ref{lem:orient}, \ref{lem:tt3} and~\ref{lem:alpha} and the degrees.

\section{How to verify}\label{sec:verify}

{\footnotesize\begin{verbatim}
42f6b3f624907602821e42c57601e311d9f3941bd6af1046e7903b11455801f4  certificates/k73_N38.json
73b58fcce089c13a090c21727106175d1424cfdbe5ff386b725481664920d840  certificates/k83_N46.json
3c8dc5150b6efd58093aa4724f14d78f9608a7afac49b0099b377e273a9911c3  certificates/k36_N72.json
\end{verbatim}}

Each file is JSON \verb|{"N": N, "a": n, "b": m, "arcs": [[u, v], ...]}| with arc $u\to v$ on
vertices $0,\dots,N-1$ (the format of \cite{PartG}), plus the construction data of
Section~\ref{sec:graphs}. Run from the folder \texttt{k36/} of the repository (on Windows,
\texttt{py} in place of \texttt{python3}); \texttt{..} is the repository root, which holds Part~G's
checker \texttt{k34-scripts/verify\_witness.py}, unchanged since its release in v0.7.0:

{\footnotesize\begin{verbatim}
python3 scripts/verify_bounds.py certificates/k73_N38.json certificates/k83_N46.json
    certificates/k36_N72.json
python3 scripts/check_constructions.py
python3 scripts/test_check_constructions.py
python3 scripts/check_same_graphs.py
python3 ../k34-scripts/verify_witness.py certificates/k73_N38.json
python3 ../k34-scripts/verify_witness.py certificates/k83_N46.json
python3 ../k34-scripts/verify_witness.py certificates/k36_N72.json
\end{verbatim}}

\noindent (The first command is one line.) Every command prints PASS and exits with status $0$;
\texttt{verify\_bounds.py} prints \verb|alpha=6 ... largest TT=2|, \verb|alpha=7 ... largest TT=2|
and \verb|alpha=2 ... largest TT=5|, then \verb|k(7,3) = r(I_7, L_3) >= 39|,
\verb|k(8,3) = r(I_8, L_3) >= 47| and \verb|k(3,6) = r(I_3, L_6) >= 73|. Times on 2026-11-11 (arm64
macOS laptop running other jobs, Python 3.15.2, low priority): \texttt{verify\_bounds.py} under
$0.1$\,s per file, the two construction scripts $0.02$\,s and $0.24$\,s, Part~G's checker $7.4$\,s,
$135.3$\,s and $115.5$\,s; \texttt{check\_same\_graphs.py} (2026-11-25, the same laptop) under
$0.1$\,s.

All five scripts are standard-library Python, make no network calls and write nothing.
\texttt{verify\_bounds.py} rejects loops, repeated arcs and $2$-cycles and computes $\alpha$ and
$\tau$ exactly by bitmask branch-and-bound; a file passes iff $\alpha\leq n-1$ and $\tau\leq m-1$.
Part~G's \texttt{verify\_witness.py} tests every $n$-subset for independence and every $m$-subset
for a transitive tournament (score sequence $0,1,\dots,m-1$).
\texttt{test\_check\_constructions.py} shows that \texttt{check\_constructions.py} rejects nine
corrupted copies of the files and that Lemmas~\ref{lem:tt3} and~\ref{lem:alpha} agree with brute
force on $600$ random small cases each. \texttt{check\_same\_graphs.py} rebuilds the two
bicirculants of \cite{Orden2} from the connection sets printed in its README and checks that the
relabellings of Section~\ref{sec:credit} carry $G_{38}$ and $G_{46}$ onto them.

\section{Published values and credit}\label{sec:credit}

\paragraph{\cite{Orden2}} T.~S.~Orden's second Zenodo dataset on this problem (published
2026-11-20, version 1.0; data and text CC~BY~4.0, code MIT) states $k(7,3)\geq 39$,
$k(8,3)\geq 47$, $k(6,4)\geq 78$ and $k(3,6)\geq 71$. Its graphs are bicirculants on
$\Z_{19}\times\{0,1\}$ and $\Z_{23}\times\{0,1\}$, a circulant on $\Z_{77}$, and a circulant on
$\Z_{70}$ with a $23$-element connection set; a second circulant, on $\Z_{69}$, gives
$k(3,6)\geq 70$. Its README says that the graphs were found by computer searches run by AI agents
under the author's direction, and that the record has not been peer reviewed. Erd\H{o}s Problem a
Day rebuilt the four graphs from their connection sets and checked them with its own program on
2026-11-22 \cite{EPAD}. On 2026-11-25 \texttt{verify\_bounds.py} accepted all five witness graphs of
the record. The bounds $k(7,3)\geq 39$ and $k(8,3)\geq 47$ are Orden's: his record is the earlier
publication.

\paragraph{The same two graphs} $G_{38}$ is the $38$-vertex graph of \cite{Orden2} relabelled, and
$G_{46}$ is its $46$-vertex graph relabelled. Both records number the vertex with residue $i$ in
orbit $a$ as $qa+i$, with $q=19$ or $23$. The map $(i,a)\mapsto(14i+16a \bmod 19,\,1-a)$ carries
the arc list of \texttt{k73\_N38.json} onto his $38$-vertex graph, and
$(i,a)\mapsto(19i+18a \bmod 23,\,1-a)$ carries the arc list of \texttt{k83\_N46.json} onto his
$46$-vertex graph (\texttt{scripts/check\_same\_graphs.py}).

\paragraph{\cite{rsv}} The GitHub repository rsvaultpudgy-create/erdos-112-lower-bounds (CC0,
created 2026-11-11 13:54 UTC, one commit, d305711817) states $k(7,3)\geq 37$ (a circulant on
$\Z_{36}$ with connection set $\{6,8,10,11,13,15\}$), $k(8,3)\geq 46$ ($\Z_{45}$,
$\{7,13,15,19,23,25,31\}$) and $k(3,6)\geq 68$ ($\Z_{67}$, a $22$-element set), and bounds for seven
other pairs. Its README says the work was produced by an AI system and is unrefereed. On 2026-11-11
\texttt{verify\_bounds.py} accepted its graphs for our three quantities.

\paragraph{\cite{Orden}} T.~S.~Orden's first Zenodo dataset on this problem (published 2026-11-09,
version 2 on 2026-11-16; data and text CC~BY~4.0, code MIT) gives a $29$-vertex witness for
$k(6,3)\geq 30$ and a written proof of $k(6,3)\leq 32$ that, it says, has not been peer reviewed. A
separate record, \cite{OrdenLean}, gives a Lean~4 proof of both bounds; it has not been built here.
\cite{Orden} states no value for the three quantities of this note, and neither does the README of
\cite{OrdenLean}.

\paragraph{\cite{IRW21}} It gives $k(4,3)=15$, $k(5,3)=23$ and $k(m,3)\leq m^2-m+3$, so
$k(7,3)\leq 45$ and $k(8,3)\leq 59$. Its closing section bounds $r(I_m,L_n)$ by an explicit
$v(m,n)$, and $v(3,6)=161$. Its recursion
$r(I_{m+1},L_{n+1})\leq 2\,r(I_{m+1},L_n)+r(I_m,L_{n+1})-1$, with $k(3,4)=21$ \cite{PartG} and the
values $r(I_2,L_5)=14$ and $r(I_2,L_6)=28$ listed in \cite{IRW21}, gives $k(3,5)\leq 55$ and
$k(3,6)\leq 137$, as \cite{rsv} also states. The certificates of \cite{PartG} were rebuilt and
replayed with a formally verified checker by Erd\H{o}s Problem a Day on 2026-11-18 \cite{EPAD2}.
\cite{Ber74} gives $k(3,3)=9$; its added-in-proof note says this result (its Proposition~2.5) also
appears in Harary and Hell, Bull.\ London Math.\ Soc.\ 6 (1974) 175--182.

\medskip
The row $k(n,3)$ as of 2026-11-25 reads $9$, $15$, $23$ for $n=3,4,5$; then $30\leq k(6,3)\leq 32$
(\cite{Orden}; \cite{rsv} also gives $30$); $39\leq k(7,3)\leq 45$; $47\leq k(8,3)\leq 59$ (lower
bounds from \cite{Orden2}).

\section{Prior-art status and what is open}\label{sec:prior}

\noindent\textbf{Not found in the sources read, as of 2026-11-25.} None of the sources read in full
or scanned (the first two lists below) states $k(3,6)\geq 73$ or a larger lower bound for $k(3,6)$;
the largest stated is $71$ \cite{Orden2}. The date 2026-11-25, here and for the row $k(n,3)$ in
Section~\ref{sec:credit}, does not cover the problem page \cite{Blo} and its forum thread, which
were last read on 2026-11-11, or the records that the list \emph{Not read} below gives as not read
again on 2026-11-25. For $k(7,3)$ and $k(8,3)$ the largest stated are $39$ and $47$
\cite{Orden2}, the values of $G_{38}$ and $G_{46}$. The assemblies of published graphs that were
built and checked on 2026-11-11 (disjoint unions and blow-ups of the graphs published by then)
reach at most $37$, $46$ and $68$. Whether one of the three graphs is an induced subgraph of a
larger published graph (for example those of McCarthy and Springfield, arXiv:2311.02135, or of
McCarthy and Monico, arXiv:2408.04067) was checked only in part, by computer runs without proof
objects, so that question stays open.

\medskip
\noindent\emph{Read in full} (prior-art reads of 2026-10-14 to 2026-10-20, of 2026-11-11, of
2026-11-17 and of 2026-11-25). Problem records: erdosproblems.com/112 with its LaTeX, history and
forum pages (page and forum again on 2026-11-11); the erdosproblems.ai entry for \#112;
teorth/erdosproblems; plasma-ai/erdos E0112; Erd\H{o}s Problem a Day (\#112 report, its
$k(3,4)=21$ page, progress page) with the two \#112 folders in pw/erdosproblemaday-verification
on 2026-11-11 (n21-attack and order21-casesplit), and on 2026-11-25 its problem page,
\cite{EPAD} with the review of \cite{Orden2}, \cite{EPAD2}, and the README of each of the three
\#112 folders added to that repository since 2026-11-11 (orden-k63, orden-small-lower-bounds and
kirtchakov-certificates-cake-lpr; their other files were not read); the forum post of 2026-09-22 by
Muhamadiev Faridun and gitlab.com/faridunislom/math.erdos112. Papers: \cite{IRW21} (arXiv v1 and
v3), \cite{Ber74}, Larson and Mitchell 1997, Erd\H{o}s and Rado 1956 and \cite{ErRa67}, Benhocine
1978, Weinert arXiv:2611.08079, Klausner and Weinert arXiv:1810.13316, McCarthy and Springfield
arXiv:2311.02135, McCarthy and Monico arXiv:2408.04067 (and the journal version), Neiman, Mackey and
Heule arXiv:2011.00683, Lohn, Lambert and Heule (FMCAD 2022), Momihara and Suda arXiv:1605.02469,
Nagy arXiv:1401.6823, Fox, He and Wigderson arXiv:2105.02383, Grzesik et al.\ arXiv:2305.19959,
Caicedo and Hilton arXiv:1510.00078, Chung 1997, Erd\H{o}s and Moser 1964 (two papers), Stearns 1959,
Moon 1971, Haddad and Sabbagh 1969 (three notes), the 1975 corrections of Harary and Hell, and the
theses of Weinert, S\'anchez-Flores and Araujo-Pardo; Reid and Parker 1970 and Choudum and
Ponnusamy 1999 (both read in full on 2026-11-11 from the publishers' open-archive PDFs);
Manoussakis and Tuza 2001 (read in full on 2026-11-17 from the publisher's open-archive PDF);
X.~Zeng, A Computer-Assisted Proof of the Oriented Ramsey Value $k(3,4)=21$ (nine pages; the PDF in
the GitHub repository ArtificialZeng/AI-Has-Taste, commit 2537a78914 of 2026-09-13, read in full on
2026-11-25; it states nothing on the three quantities of this note). Data: \cite{Orden} (all seven
files of version 1; of version 2, the README and the proof, its other five files being identical to
version 1's); \cite{Orden2} (the README, the definitions, the three scripts and their output; the
five arc lists checked by program); the README of \cite{OrdenLean}; the README of
Tyorden/erdos-112-small-bounds-lean; the file math/directed-ramsey/results/k112-2026-11-19.md of
the GitHub repository adineo-cmd/open-research-lab, which states $k(4,4)\geq 40$ and
$k(3,5)\geq 39$ from circulants and nothing on the three quantities of this note; \cite{rsv}
(README, log, witness lists and logs); OEIS A122026, A122027, A224886, A003141, A001225, A362137
and A027688, with W.~D.~Smith's pages linked from A122026; Certify Parts G and J.

\medskip
\noindent\emph{Scanned, not read in full:} arXiv math.CO from 2026-10-19 to 2026-11-10 ($1{,}122$
titles and abstracts) and $16$ targeted arXiv queries; $29$ DataCite queries; $10$ GitHub
searches; the works citing \cite{IRW21} in OpenAlex, Semantic Scholar, OpenCitations, Crossref and
zbMATH (two: Klausner and Weinert, and \cite{Orden}); arXiv:2604.23433, 2608.12207 and 2403.02298
in part. On 2026-11-25, for the window since 2026-11-11: six arXiv queries, eight Zenodo and
DataCite queries, $44$ GitHub searches (repositories, commits, issues and code) with the READMEs of
$74$ of the repositories found (five of them cut at $512{,}000$ bytes), two web searches, and the
works citing \cite{IRW21} in OpenAlex (four: the two above, \cite{Orden2} and an SSRN record with
the title of X.~Zeng's paper listed above, doi:10.2139/ssrn.7416699, which was not opened); also the
description of the Zenodo record of J.~Mackey and A.~T.~Mat, The directed Ramsey number $R(7)$ is at
most $45$ (doi:10.5281/zenodo.24446589, 2026-11-22), which concerns tournaments; and four searches
of the public data of the Palomar registry of Lean~4 formalizations (palomar-registry.org) and a
search of its feed of $200$ recent entries, which returned two entries on \#112, both Orden's:
PALOMAR-2026-11-15-000004 (his bounds on $k(6,3)$) and PALOMAR-2026-11-22-000008 (the four bounds
of \cite{Orden2}).

\medskip
\noindent\emph{Not read:}
\begin{itemize}
\item The Lean files of \cite{OrdenLean} and of Tyorden/erdos-112-small-bounds-lean, whose README
says they prove the four bounds of \cite{Orden2}. Nothing of either was built here.
\item Not read again since 2026-11-11: erdosproblems.com/112 and its forum thread. \cite{Orden2}
reports a read of the thread on 2026-11-20 that found no bound above its own for its four cases.
\item Not read again on 2026-11-25: the OEIS entries and the $k(3,4)=21$ page of Erd\H{o}s Problem a
Day. Its \#112 report and its progress page were fetched on 2026-11-25 and searched by program for
values of the three quantities, and were not read in full again.
\item The files of the three \#112 folders added to pw/erdosproblemaday-verification since
2026-11-11, apart from their READMEs. \cite{EPAD} and \cite{EPAD2} name these folders as the
bundles of their checks.
\item M.~Chen, Y.~Kou, X.~Ma and S.~Zhang, TT3-factors in oriented graphs with low independence
number, Discrete Math.\ 349 (2026) 115229, doi:10.1016/j.disc.2026.115229.
\item Library only: F.~Harary and P.~Hell, Generalized Ramsey theory for graphs V. The Ramsey number
of a digraph, Bull.\ London Math.\ Soc.\ 6 (1974) 175--182 (main text); J.~Bang-Jensen and
F.~Havet, Tournaments and semicomplete digraphs, chapter 2 of \emph{Classes of Directed Graphs}
(Springer, 2018), section 2.7.1; A.~Benhocine's 1977 thesis (Paris VI).
\item Not reachable from our machine: the Google Scholar and MathSciNet lists of works citing
\cite{IRW21}.
\item Also not read in full: the ordinal-partition sources of Rosenstein (1982, chapter 11),
Williams (1977, chapter 7), Nosal, Milner, and Hajnal and Larson (2010); Wang, Wang and Yan,
arXiv:2606.14161; Chen, Miao and Zhou, J.\ Comb.\ Optim.\ 51 (2026) 15 (publisher preview only).
\end{itemize}

\noindent The exact values of all three numbers are open.

\section{What would falsify this}\label{sec:falsify}

\begin{itemize}
\item Seven pairwise non-adjacent vertices or a transitive triangle in \texttt{k73\_N38.json}; eight
pairwise non-adjacent vertices or a transitive triangle in \texttt{k83\_N46.json}; three pairwise
non-adjacent vertices or a transitive tournament on six vertices in \texttt{k36\_N72.json}. Part~G's
checker prints such a set when it meets one.
\item A file whose SHA-256 differs from Section~\ref{sec:verify}.
\item For priority, not for truth: a source dated before the release that states $k(3,6)\geq 73$ or
a larger lower bound for $k(3,6)$.
\item For the entry of Section~\ref{sec:credit} on $G_{38}$ and $G_{46}$, which is not a claim to a
bound: a failure of \texttt{scripts/check\_same\_graphs.py}.
\end{itemize}

\section{Questions}\label{sec:questions}

\begin{question}
The three numbers lie in $[39,45]$, $[47,59]$ and $[73,137]$ (the last using $k(3,4)=21$). Where
are the true values? The bound $m^2-m+3$ of \cite{IRW21} is attained for $m=3$, $4$ and $5$, and
\cite{Orden} puts $k(6,3)$ below it (a written proof that, by its own account, is not peer
reviewed, and a Lean~4 proof in \cite{OrdenLean} that was not built here); how far below it are
$k(7,3)$ and $k(8,3)$?
\end{question}

\begin{question}
$G_{38}$ and $G_{46}$ have a semiregular automorphism of prime order with two orbits, and $G_{72}$
is a Cayley digraph of a non-cyclic abelian group. Is there an algebraic family behind them that
reaches other pairs $(n,m)$?
\end{question}

\begin{question}
Two records, \cite{Orden2} and the files of this release, hold the same graph on $38$ vertices and
the same graph on $46$ vertices, up to relabelling. Is $G_{38}$ the only oriented graph on $38$
vertices with no transitive triangle and no seven independent vertices that has an automorphism of
order $19$ with two orbits, and $G_{46}$ the only one of its kind on $46$ vertices? What happens
with no symmetry assumed?
\end{question}

\begin{thebibliography}{OrdenLean}

\bibitem[Ber74]{Ber74} J.-C. Bermond, \emph{Some Ramsey numbers for directed graphs}, Discrete
Math.\ 9 (1974) 313--321, doi:10.1016/0012-365X(74)90077-6.

\bibitem[Blo]{Blo} T.~F.~Bloom, \emph{Erd\H{o}s Problem \#112},
\url{https://www.erdosproblems.com/112}, accessed 2026-11-11.

\bibitem[EPAD]{EPAD} Erd\H{o}s Problem a Day (P.~White with Claude), \emph{\#112, oriented Ramsey
numbers: $k(6,3)$ is 30, 31 or 32. Tyler Orden's bounds check out, and four more small cases move
up}, 2026-11-22,
\url{https://erdosproblemaday.com/day/112-orden-k63-between-30-and-32-and-four-new-lower-bounds},
with the review of \cite{Orden2} at \url{https://erdosproblemaday.com/submission/eqf2h6}; accessed
2026-11-25.

\bibitem[EPAD2]{EPAD2} Erd\H{o}s Problem a Day, page of 2026-11-18 on the certificates of
\cite{PartG},
\url{https://erdosproblemaday.com/day/112-kirtchakov-k34-certificates-checked-by-a-verified-checker};
accessed 2026-11-25.

\bibitem[ErRa67]{ErRa67} P.~Erd\H{o}s and R.~Rado, \emph{Partition relations and transitivity
domains of binary relations}, J.\ London Math.\ Soc.\ 42 (1967) 624--633,
doi:10.1112/jlms/s1-42.1.624.

\bibitem[IRW21]{IRW21} F.~Ihringer, D.~Rajendraprasad and T.~Weinert, \emph{New bounds on the
Ramsey number $r(I_m,L_n)$}, Discrete Math.\ 344 (2021), no.~3, 112268; arXiv:1707.09556.

\bibitem[Orden]{Orden} T.~S.~Orden, \emph{New bounds on the Erd\H{o}s-Rado number $k(6,3)$:
$30\leq k(6,3)\leq 32$}, Zenodo (2026-11-09),
\href{https://doi.org/10.5281/zenodo.24064871}{doi:10.5281/zenodo.24064871}; version 2
(2026-11-16), doi:10.5281/zenodo.24272287.

\bibitem[Orden2]{Orden2} T.~S.~Orden, \emph{New lower bounds for small Erd\H{o}s-Rado numbers
$k(7,3)$, $k(8,3)$, $k(6,4)$ and $k(3,6)$ (Erd\H{o}s problem \#112)}, Zenodo (2026-11-20), version
1.0, \href{https://doi.org/10.5281/zenodo.24390891}{doi:10.5281/zenodo.24390891}.

\bibitem[OrdenLean]{OrdenLean} T.~S.~Orden, \emph{Erd\H{o}s Problem \#112: a machine-checked proof
that $30\leq k(6,3)\leq 32$ (Lean~4)}, Zenodo, version 1.3 (2026-11-16),
doi:10.5281/zenodo.24272901; all versions, doi:10.5281/zenodo.24237180.

\bibitem[PartG]{PartG} D.~Kirtchakov, \emph{An Erd\H{o}s-Rado oriented Ramsey number determined:
$k(3,4)=r(I_3,L_4)=21$}, Certify v0.7.0 (2026), doi:10.5281/zenodo.21890619.

\bibitem[rsv]{rsv} rsvaultpudgy-create, \emph{erdos-112-lower-bounds}, GitHub, CC0, commit
d305711817 (2026-11-11), \url{https://github.com/rsvaultpudgy-create/erdos-112-lower-bounds}.

\end{thebibliography}

\end{document}
